\documentclass[10.5pt]{article}
\usepackage[utf8]{inputenc}
\usepackage[margin=0.75in]{geometry}
\usepackage{setspace}
\setstretch{1.12}
\usepackage{hyperref}
\hypersetup{colorlinks=true, linkcolor=black, urlcolor=blue}

\begin{document}
\noindent \textbf{Danil Popov} \\
Independent Researcher \\
E-mail: \href{mailto:perque.balom@gmail.com}{perque.balom@gmail.com} \\
ORCID: \href{https://orcid.org/0009-0008-4054-122X}{0009-0008-4054-122X} \\
Replication Archive: \url{https://github.com/karodg/compute-grid-collision} \\
\vspace{0.3cm}

\noindent The Editors \\
\textit{Energy Policy} \\
Elsevier \\
\vspace{0.25cm}

\noindent \textbf{Subject: Submission of Research Manuscript --- `The Compute-Grid Collision: AI Load Shocks, Wholesale Capacity Rents, and the Regressive Ratepayer Tax''} \\
\vspace{0.2cm}

\noindent Dear Editors, \\
\vspace{0.15cm}

\noindent I am pleased to submit my original research manuscript entitled \textbf{`The Compute-Grid Collision: AI Load Shocks, Wholesale Capacity Rents, and the Regressive Ratepayer Tax''} for consideration for publication in \textit{Energy Policy}.

The rapid build-out of artificial intelligence infrastructure requires continuous baseload electric power, driving technology hyperscalers into behind-the-meter co-location agreements and multi-gigawatt interconnection queues. This paper provides the first general-equilibrium and quasi-experimental econometric investigation into how this computational shock disrupts power grid economics and utility ratemaking. Because hyperscalers deploy specialized GPU clusters with multi-billion-dollar capital costs, idle costs dominate input energy prices by several orders of magnitude, rendering hyperscaler electricity demand hyper-inelastic. When this inelastic load shock collides with steep wholesale capacity supply curves and transmission bottlenecks, capacity clearing prices spike to administrative caps, which regulated utilities pass directly into volumetric residential tariffs via formula-rate mechanisms (FERC Attachment H and retail SOS riders).

Exploiting the unprecedented July 2024 PJM Base Residual Auction (BRA) clearing price shock alongside the landmark November 1, 2024 FERC rejection of the Talen--Amazon Susquehanna co-location agreement (Docket ER24-2172), the paper establishes four core empirical and theoretical contributions:
\begin{itemize}
    \item \textbf{Econometric Pass-Through \& 2SLS IV Identification:} Across a balanced panel of 120 PJM counties (2019Q1--2026Q2,  = 3,600$), two-way fixed effects regressions demonstrate that wholesale capacity price spikes pass through into retail residential bills with an elasticity of 0.38 (0.77 cents/kWh per \/MW-day,  < 0.01$). Instrumenting data center load share with distance to 500 kV transmission substations interacted with national GPU deployment shocks confirms causality (First-stage  = 78.42$, $\beta_{IV} = 3.821***$).
    \item \textbf{Spatial Boundary Discontinuity Design (RDD):} Boundary RDD across contiguous county borders separating constrained Locational Deliverability Areas (BGE, DOM) from unconstrained zones reveals an immediate, discontinuous retail tariff jump of \textbf{+3.214 cents/kWh} ( < 0.001$) directly at the administrative boundary.
    \item \textbf{Industrial Manufacturing Crowd-Out:} Energy-intensive, trade-exposed manufacturing (NAICS 325 Chemical and NAICS 331 Primary Metal) suffers an immediate \textbf{18.4\% employment contraction} ($-420$ jobs per border county,  < 0.001$), while event studies confirm flat parallel pre-trends ( = 0.38, p = 0.912$) before declining to a cumulative $-21.8\%$ contraction by quarter  = +7$.
    \item \textbf{Regressive Distributional Incidence \& Counterfactuals:} Bottom-quintile households incur an energy burden expansion of \textbf{+2.38 percentage points} (+25.7\% relative expansion), compared to only +0.21 pp (+14.8\%) for top-quintile earners (Regressivity Multiplier = 1.74x). I simulate a Beneficiary-Pays Incremental Tariff (BPIT) reform assigning network capex to hyperscalers, reducing residential tariff pass-through by 82\%.
\end{itemize}

This manuscript represents original work, has not been published previously, and is not currently under consideration elsewhere. Thank you very much for your time and editorial evaluation. \\
\vspace{0.25cm}

\noindent Sincerely, \\
\textbf{Danil Popov} \\
Independent Researcher \\
perque.balom@gmail.com

\end{document}

